\uuid{ocOA}
\titre{Arithmétique}
\auteur{bodin}
\organisation{exo7}
\licence{CC BY-NC-SA 4.0}
\datecreate{2020-01-21}
\niveau{L1}
\module{Algèbre}
\chapitre{Arithmétique dans Z}
\sousChapitre{Autre}
\difficulte{}
\theme{Arithmétique}
\exo7qcmid{arnaud.arnaud.019}

\contenu{
\begin{question}
\qclassification{103.99}
Quelles sont les assertions vraies ?
\begin{answers}
    \good{Soit $P\in \Cc[X]$. $X-a$ divise $P$ ssi $P(a)=0$.}

    \good{Soit $P\in \Rr[X]$ de degré impair. Il existe $x \in \Rr$ tel que $P(x) = 0$.}

    \good{Soit $P\in \Rr[X]$, les racines de $P^2$ sont d'ordre au moins $2$.}

    \bad{Soit $P\in \Rr[X]$, $x$ est racine simple ssi $P(x) = 0$.}

    \bad{Un polynôme  $P\in \Cc[X]$ de degré $n$ a $n$ racines réelles.}

\end{answers}
\end{question}
}
